\section{测地线的初值问题}
\label{sec:riem-geodesic-ivp}

在平面上，“一直向前走”就是沿直线匀速运动，加速度为零。在球面上照搬这个说法会出问题：从$p=(1,0,0)$沿切方向$v=(0,1,0)$走直线$p+tv$，长度$\sqrt{1+t^2}$大于$1$，立刻离开球面。沿赤道$\gamma(t)=(\cos t,\sin t,0)$匀速前进则始终留在球面上，它的环境加速度$\ddot\gamma=-\gamma$不为零，却完全指向球面法向，只是“贴住球面”所必需的。沿球面没有转弯，意味着加速度的切向部分为零，也就是速度的协变导数为零。

\begin{definition}[测地线]
\label{def:riem-geodesic}
黎曼流形$(M,g)$上的光滑曲线$\gamma:I\to M$（$I$为区间）若满足
\[
\frac{D\dot\gamma}{\mathrm dt}=0,
\]
其中协变导数取Levi-Civita联络，就称$\gamma$为\term{测地线}（geodesic）。换言之，测地线的速度沿它自身平行移动。
\end{definition}

由定理~\ref{thm:riem-transport-isometry}，速度的长度沿测地线不变，所以测地线总是匀速的。这也意味着“测地线”是带参数化的概念：赤道换成非匀速参数化$t\mapsto(\cos t^3,\sin t^3,0)$，像集不变，速度大小却在变化，因此不再满足定义；只有仿射换元$t=as+b$保持测地线方程。本节先研究由初始位置和初始速度确定测地线的微分方程，再用指数映射把从一点出发的所有测地线汇总起来。

\subsection{测地线方程}
\label{sec:riem-geodesic-equation}

在坐标图中写$\gamma(t)$的坐标为$x(t)$，则$\dot\gamma=\sum_j\dot x^j\partial_j$。在式\eqref{eq:riem-covariant-derivative-along-curve}中取$V=\dot\gamma$，即$V^j=\dot x^j$，测地线条件化为
\begin{equation}
\ddot x^k(t)+\sum_{i,j=1}^n\Gamma^k_{i,j}\bigl(x(t)\bigr)\,\dot x^i(t)\,\dot x^j(t)=0,\qquad k=1,\ldots,n .
\label{eq:riem-geodesic-equation}
\end{equation}
与平行移动方程\eqref{eq:riem-parallel-equation}不同，这里系数依赖未知的$x(t)$，方程关于速度是二次的，因此是非线性方程。

\begin{example}[极坐标中的直线与球面上的大圆]
\label{ex:riem-geodesic-examples}
欧氏平面的极坐标中，由例~\ref{ex:riem-christoffel-examples}，测地线方程为
\[
\ddot r-r\dot\theta^2=0,\qquad\ddot\theta+\frac2r\dot r\dot\theta=0 .
\]
直线$x=1$匀速参数化为$y=t$，对应$r=\sqrt{1+t^2}$、$\theta=\arctan t$。计算得$\dot r=t(1+t^2)^{-1/2}$、$\ddot r=(1+t^2)^{-3/2}$、$\dot\theta=(1+t^2)^{-1}$、$\ddot\theta=-2t(1+t^2)^{-2}$，代入两式均为零。直线在直角坐标中的方程$\ddot x=\ddot y=0$，在极坐标中出现了Christoffel项。

单位球面上，方程为
\[
\ddot\theta-\sin\theta\cos\theta\,\dot\phi^2=0,\qquad\ddot\phi+2\cot\theta\,\dot\theta\dot\phi=0 .
\]
赤道$\theta=\pi/2$、$\phi=t$满足两式，是测地线。纬线$\theta=\theta_0\neq\pi/2$、$\phi=t$使第一式左端等于$-\sin\theta_0\cos\theta_0\neq0$，所以不是测地线：沿北半球的纬线匀速前进，需要一直向北侧转弯。
\end{example}

\paragraph{局部存在唯一与初值依赖}
把二阶方程改写成关于$(x,\dot x)$的一阶系统$\dot x=w$、$\dot w^k=-\sum_{i,j}\Gamma^k_{i,j}(x)w^iw^j$，右端是$(x,w)$的光滑函数。对这类方程，下面的结论成立，证明同样归入前面所引的卷六相应章节。

\begin{theorem}[光滑常微分方程的局部理论]
\label{thm:riem-smooth-ode}
设$W\subseteq\mathbb R^m$开，$F:W\to\mathbb R^m$光滑。则对每个$y_0\in W$：
\begin{enumerate}
\item 存在$\varepsilon>0$及唯一的解$y:(-\varepsilon,\varepsilon)\to W$，满足$\dot y=F(y)$、$y(0)=y_0$；
\item 两个解在公共定义区间上若有一个时刻取值相同，则在整个公共区间上相同，因而所有解可以拼成唯一的最大解，其定义区间是包含$0$的开区间；
\item 最大解$y(t;y_0)$关于$(t,y_0)$光滑，所有使其有定义的$(t,y_0)$构成开集；
\item 若最大解的定义区间右端点$t^+$有限，则当$t\to t^+$时$y(t)$终将离开$W$的每个紧子集。
\end{enumerate}
\end{theorem}

应用到测地线方程，得到：对每个$p\in M$和$v\in T_pM$，存在唯一的最大测地线$\gamma_v:I_v\to M$，满足$\gamma_v(0)=p$、$\dot\gamma_v(0)=v$，其中$I_v$是包含$0$的开区间。跨图时的拼接与平行移动相同，由唯一性保证一致。光滑依赖说明，初始速度稍有变化，测地线在有限时间内的位置也只稍有变化。

\paragraph{初始速度与时间的换算}
初始速度放大$c$倍，等价于把时间加快$c$倍：对$c\neq0$，$t\mapsto\gamma_v(ct)$满足测地线方程，初值为$p$与$cv$，所以
\begin{equation}
\gamma_{cv}(t)=\gamma_v(ct),\qquad t\in I_{cv}=\tfrac1cI_v .
\label{eq:riem-geodesic-rescaling}
\end{equation}
验证只需注意方程\eqref{eq:riem-geodesic-equation}的每一项关于时间导数是二次齐次的：$\ddot x$与$\dot x^i\dot x^j$都乘上$c^2$。

\paragraph{局部解未必能延拓到所有时间}
定理~\ref{thm:riem-smooth-ode}只保证解在某个时间区间内存在。在开单位球$M=\{x\in\mathbb R^n:\lVert x\rVert_2<1\}$上取欧氏度量，测地线是匀速直线。从原点以速度$2\bm e_1$出发，$\gamma(t)=2t\bm e_1$只在$t\in(-1/2,1/2)$内留在$M$中，最大定义区间是有限的；这正是结论(4)所说的“离开每个紧子集”。完整的单位球面上则不同，大圆对所有实数时间都有定义。测地线能否延拓到所有时间，取决于流形的整体性质，第~\ref{sec:riem-shortest-paths}节的Hopf--Rinow定理回答这个问题。

\subsection{黎曼指数映射}
\label{sec:riem-exponential-map}

从一点$p$出发，每个初始速度$v$给出一条测地线。把“沿$v$方向走单位时间到达哪里”记录下来，就得到从切空间$T_pM$到流形的一个映射。它把$p$附近的点用一个切向量来标记，类似于在地图上用“方向和距离”标记位置。

\begin{definition}[黎曼指数映射]
\label{def:riem-exponential-map}
对$p\in M$，令$\mathcal E_p=\{v\in T_pM:1\in I_v\}$为测地线$\gamma_v$至少存在到时刻$1$的初速度全体。定义
\[
\exp_p:\mathcal E_p\to M,\qquad\exp_p(v)=\gamma_v(1),
\]
称为$p$处的\term{黎曼指数映射}（Riemannian exponential map）。
\end{definition}

由式\eqref{eq:riem-geodesic-rescaling}，若$v\in\mathcal E_p$且$0\leq t\leq1$，则$tv\in\mathcal E_p$，并且
\[
\exp_p(tv)=\gamma_{tv}(1)=\gamma_v(t).
\]
所以沿射线$t\mapsto tv$的像正是测地线$\gamma_v$，定义域$\mathcal E_p$关于原点呈星形。由定理~\ref{thm:riem-smooth-ode}的结论(3)，$\mathcal E_p$是包含$0$的开集，$\exp_p$在其上光滑。

\begin{theorem}[指数映射的局部可逆性]
\label{thm:riem-exp-local-diffeo}
对每个$p\in M$，$d(\exp_p)_0:T_0(T_pM)\cong T_pM\to T_pM$是恒等映射。因此存在$0$在$T_pM$中的开邻域$V$和$p$在$M$中的开邻域$U$，使$\exp_p:V\to U$是微分同胚。
\end{theorem}

\begin{proof}
线性空间$T_pM$在$0$处的切空间可以认同为$T_pM$本身。对$v\in T_pM$，曲线$t\mapsto tv$在$t=0$的速度是$v$，它在$\exp_p$下的像是$\gamma_v(t)$，速度为$\dot\gamma_v(0)=v$。由链式法则，$d(\exp_p)_0(v)=v$。微分在一点可逆，由反函数定理即得局部微分同胚；反函数定理的欧氏形式可参考\BookChapterRef{chap:implicit-inverse-functions}{卷二“隐函数与反函数”章}，它在流形上的形式见第~\ref{chap:manifold-relations}章。
\end{proof}

若$V$还可以取为关于原点星形的，例如以$0$为心的球$\{v:\lVert v\rVert_g<\varepsilon\}$，则称像集$U=\exp_p(V)$为$p$的\term{正常邻域}（normal neighborhood）。在正常邻域中，每个点$q$都唯一地写成$q=\exp_p(v)$，$v\in V$，从$p$到$q$有一条“径向”测地线$t\mapsto\exp_p(tv)$。取$T_pM$的一组标准正交基把$v$写成坐标，就得到$U$上的一组坐标，称为$p$处的法坐标。

\begin{example}[单位球面的指数映射]
\label{ex:riem-sphere-exp}
单位球面$\mathbb S^{n}\subseteq\mathbb R^{n+1}$上，取$p$与$v\perp p$，$r=\lVert v\rVert_2>0$。曲线$\gamma(t)=\cos(tr)\,p+\sin(tr)\,v/r$满足$\lVert\gamma(t)\rVert_2=1$，$\gamma(0)=p$，$\dot\gamma(0)=v$，且$\ddot\gamma=-r^2\gamma$为纯法向，所以它是测地线。于是$\mathcal E_p=T_p\mathbb S^n$，并且
\[
\exp_p(v)=\cos\lVert v\rVert_2\,p+\frac{\sin\lVert v\rVert_2}{\lVert v\rVert_2}\,v,\qquad\exp_p(0)=p .
\]
所有长度为$\pi$的初速度都到达对跖点$-p$，所以$\exp_p$不是单射；在开球$\lVert v\rVert_2<\pi$上它是到$\mathbb S^n\setminus\{-p\}$的微分同胚。可见即使指数映射处处有定义，可逆性也只是局部的。
\end{example}

\paragraph{与矩阵指数的区别}
第~\ref{chap:lie-groups}章会在矩阵Lie群上用矩阵指数$e^{\bm A}=\sum_k\bm A^k/k!$把单位元处的切向量变成群元素。两者都叫“指数”，都从切空间出发、在零点处微分为恒等，但来源不同：黎曼指数映射由度量经Levi-Civita联络确定，换一个度量就会改变；矩阵指数只由群的乘法确定，不需要度量。在带特殊度量的某些群上二者可以一致，这需要单独验证，不能仅凭名称相同就互相替代。
